\chapter{Desenvolvimento}\label{cap:desenvolvimento}

Como orientado pelo SIBI \cite{ManualSibi2025}, essa é a "parte principal do texto que contém a exposição ordenada e pormenorizada do assunto. Divide-se em seções e subseções, que variam em função da abordagem do tema e do método".

O texto referente ao desenvolvimento do trabalho versa sobre como o pesquisador executou o método de pesquisa. Assim, deverá descrever os dados coletados e analisá-los conforme o método prescreve. Alguns pesquisadores preferem descrever todos os dados e, em outra seção, analisá-los. Outros preferem descrever e analisar dado a dado.

\section{Construção de artefatos}
Caso o trabalho de pesquisa demande a construção de algum artefato (tecnológico ou teórico), sugerimos que seja feita uma seção específica para descrição do desenvolvimento desse artefato antes de usar os métodos de pesquisa que terão o objetivo de coletar e analisar os dados oriundos do uso desse artefato.

\section{Coleta de dados}
É o momento em que o pesquisador descreve para o leitor os dados que coletou. Seria como uma “foto” do fenômeno que o pesquisador observou. Não cabem críticas ou análises, apenas a exposição do que foi visto.

\section{Resultados da pesquisa}
No momento da análise dos dados, o pesquisador deve associar a teoria levantada no referencial teórico com os dados descritos anteriormente. Assim, o pesquisador realça se o que verificou na prática está de acordo, ou não, com a literatura.  O pesquisador pode e deve fazer comentários analíticos quanto ao que está diferente do que prega a teoria, realçando se tal diferença foi prejudicial, ou não.

É interessante salientar, também, o que está em consonância com o referencial teórico e verificar se essa concordância surtiu os efeitos benéficos citados pela teoria, ou não.

Essas análises feitas ao longo do desenvolvimento constituem os resultados da sua pesquisa. A seção de análise é comumente denominada "Resultados da pesquisa".

\section{Discussão dos resultados}
Após a apresentação dos resultados, o pesquisador pode criar uma seção de "Discussão dos resultados", em que confronta seus resultados com trabalhos correlatos, provavelmente citados na seção de referencial teórico.

\newpage
\textbf{Observações:} 
\begin{enumerate}
	\item As seções apresentadas anteriormente podem estar em um subnível dentro de uma seção maior chamada "Desenvolvimento", Como apresentado nesse guia; ou podem estar no mesmo nível da estrutura de seções que a Introdução, referencial e Metodologia;
	\item Caso o método utilizado fosse um experimento de um determinado aplicativo, antes da descrição dos dados coletados e de sua análise, haveria o relato do desenvolvimento do artefato que está sendo experimentado. Após isso, o pesquisador faria a descrição do ambiente pré-intervenção, a descrição da intervenção em si (incluindo configuração de ambiente necessária, softwares criados ou utilizados, etc) e a descrição pós-intervenção. A seguir, o autor poderia analisar as alterações percebidas com a intervenção à luz da teoria compulsada na seção de referencial teórico.
	\item Em um estudo de caso, o pesquisador poderia descrever caso a caso e depois analisá-los em busca de categorias que expliquem o fenômeno. Perceba que, dependendo do método de pesquisa utilizado, haverá diferenças na seção de desenvolvimento.
\end{enumerate}



